\chapter{Laryngectomy}

\section*{Introduction \& Clinical Indications}
Laryngectomy is a major surgical procedure involving the partial or complete removal of the larynx, commonly known as the voice box. This complex operation is primarily performed for the treatment of advanced laryngeal cancer, but can also be indicated for severe laryngeal trauma, radiation necrosis, or other extensive laryngeal pathologies. The procedure fundamentally alters the patient's airway, separating it permanently from the digestive tract, and necessitates the creation of a permanent opening in the neck for breathing, known as a tracheostoma. Consequently, patients undergoing laryngectomy experience significant changes in voice production and swallowing mechanisms, requiring extensive rehabilitation to restore these vital functions. The extent of the resection, whether partial or total, depends on the tumor's size, location, and invasiveness, aiming for complete oncologic clearance while preserving as much function as possible.

\begin{itemize}
  \item Total Laryngectomy
  \item Partial Laryngectomy (e.g., Supraglottic Laryngectomy, Hemilaryngectomy, Supracricoid Laryngectomy)
  \item Voice Box Removal
  \item Laryngeal Resection
  \item Laryngopharyngectomy (if pharynx is also extensively involved)
  \item Tracheostoma Creation (as part of the procedure)
\end{itemize}

The fundamental principle of laryngectomy, particularly total laryngectomy, is the complete extirpation of the laryngeal tumor with clear surgical margins, adhering to established oncologic principles for head and neck cancer. This involves the en bloc removal of the entire larynx, including the thyroid cartilage, cricoid cartilage, epiglottis, and vocal cords, along with any involved surrounding tissues or lymph nodes (via concomitant neck dissection). The primary rationale is to achieve definitive local control of advanced laryngeal malignancies that are either too extensive for organ preservation strategies (e.g., radiation or chemoradiation) or have recurred after such treatments.

Functionally, the procedure aims to permanently separate the respiratory tract from the digestive tract. This is crucial for preventing aspiration of food and liquids into the lungs, which is a significant and life-threatening risk in patients with extensive laryngeal disease or after partial resections that compromise glottic closure. Following the removal of the larynx, the trachea is brought forward and sutured to an opening in the neck skin, creating a permanent tracheostoma for breathing. The pharynx, which is transected during the removal of the larynx, is then meticulously reconstructed and closed to restore continuity of the digestive pathway. For voice rehabilitation, various methods are employed, including the use of an electrolarynx, esophageal speech, or the surgical creation of a tracheoesophageal puncture (TEP) with placement of a voice prosthesis. The choice of procedure, whether total or partial, is guided by the extent of the disease, with partial laryngectomies preserving some laryngeal structures and potentially some voice function, but carrying a higher risk of aspiration.

The history of laryngectomy is a testament to surgical innovation and the relentless pursuit of improving patient outcomes. The first successful total laryngectomy was performed by Theodor Billroth in Vienna in 1873 on a 36-year-old patient with laryngeal cancer. This groundbreaking operation, though initially met with skepticism, demonstrated the feasibility of removing the entire larynx and establishing a permanent airway. Early procedures were associated with high morbidity and mortality, primarily due to infection and aspiration, and patients were left without a voice, relying on writing or gestures for communication.

Over the subsequent decades, surgical techniques were refined, and the understanding of laryngeal anatomy and pathology deepened. The early 20th century saw the development of partial laryngectomy techniques, pioneered by surgeons like Chevalier Jackson, which aimed to preserve portions of the larynx and some voice function for less extensive cancers. The introduction of reconstructive techniques and, crucially, methods for voice rehabilitation, such as the electrolarynx in the mid-20th century and later the development of tracheoesophageal puncture (TEP) and voice prostheses by Dr. Eric Blom and Dr. Mark Singer in the late 1970s, revolutionized post-laryngectomy quality of life. These advancements transformed laryngectomy from a purely ablative, life-saving procedure into one that also prioritized functional restoration and patient well-being, significantly improving communication and swallowing outcomes for patients.

\begin{itemize}
  \item To treat advanced squamous cell carcinoma of the larynx (T3 or T4 disease) that is not amenable to organ preservation strategies or has failed prior non-surgical treatments.
  \item For recurrent laryngeal cancer after definitive radiation therapy or chemoradiation, where salvage surgery is indicated.
  \item To manage severe laryngeal chondronecrosis, often a complication of radiation therapy, which causes intractable pain, airway obstruction, or persistent infection.
  \item In cases of extensive laryngeal trauma that results in irreparable damage to the laryngeal framework, vocal cords, or airway integrity.
  \item For large, invasive tumors of the hypopharynx or cervical esophagus that extend into the larynx, requiring en bloc resection.
  \item To address persistent and severe aspiration that is refractory to all other medical and surgical interventions, significantly impacting pulmonary health.
  \item As part of a multidisciplinary treatment plan for certain rare laryngeal sarcomas or other aggressive non-squamous malignancies.
  \item When a patient presents with a compromised airway due to laryngeal obstruction from a tumor, and immediate surgical intervention is required.
\end{itemize}

Laryngectomy can serve as both a primary and a secondary procedure, depending on the clinical context. It is most commonly considered a primary curative surgery for patients presenting with advanced-stage laryngeal cancer (e.g., T3 or T4 disease) where the extent of the tumor precludes organ preservation with radiation or chemoradiation, or where such non-surgical approaches are deemed unlikely to achieve local control. In these scenarios, laryngectomy is the initial definitive treatment aimed at complete tumor eradication. Conversely, laryngectomy frequently functions as a secondary or salvage procedure. This occurs when patients have undergone prior definitive radiation therapy or chemoradiation for laryngeal cancer, but experience local recurrence or persistent disease. In such cases, laryngectomy offers the best chance for cure after the failure of initial non-surgical management. Less commonly, it may be performed as a palliative measure to relieve severe symptoms like airway obstruction or intractable aspiration in patients with advanced, incurable disease, though this is typically a last resort after all other less invasive options have been exhausted.

\section*{Relevant Surgical Anatomy}
The larynx, or voice box, is a complex cartilaginous structure located in the anterior neck, superior to the trachea and inferior to the hyoid bone. It is composed of several cartilages, including the large thyroid cartilage (forming the laryngeal prominence or "Adam's apple"), the complete ring-shaped cricoid cartilage (the only complete cartilaginous ring of the airway), the leaf-shaped epiglottis (which protects the airway during swallowing), and paired arytenoid cartilages (crucial for vocal cord movement). These cartilages are connected by ligaments and membranes, such as the cricothyroid membrane and thyrohyoid membrane. The interior of the larynx is divided into three regions: the supraglottis (above the vocal cords, including the epiglottis and false vocal cords), the glottis (containing the true vocal cords and the rima glottidis), and the subglottis (below the vocal cords to the inferior border of the cricoid cartilage).

The larynx is richly supplied by the superior laryngeal artery (a branch of the superior thyroid artery, which originates from the external carotid artery) and the inferior laryngeal artery (a branch of the inferior thyroid artery, from the thyrocervical trunk). Venous drainage mirrors the arterial supply, with superior laryngeal veins draining into the superior thyroid vein (then to the internal jugular vein) and inferior laryngeal veins draining into the inferior thyroid vein (then to the brachiocephalic vein). Innervation is provided by branches of the vagus nerve (cranial nerve X): the superior laryngeal nerve (internal branch provides sensation above the vocal cords, external branch supplies the cricothyroid muscle) and the recurrent laryngeal nerve (supplies all other intrinsic laryngeal muscles and sensation below the vocal cords). Lymphatic drainage is extensive and critical in cancer staging, involving prelaryngeal (Delphian), pretracheal, paratracheal, and deep cervical lymph nodes (levels II, III, IV, and sometimes VI). Key surgical landmarks include the hyoid bone, thyroid cartilage, cricoid cartilage, and the strap muscles (sternohyoid, sternothyroid, omohyoid, thyrohyoid) which overlie the larynx and are often divided or retracted during surgery. The carotid sheath, containing the common carotid artery, internal jugular vein, and vagus nerve, lies lateral to the larynx and must be carefully preserved during neck dissection.

\section*{Preoperative Workup \& Preparation}
Thorough pre-operative preparation is crucial for patients undergoing laryngectomy to optimize outcomes and facilitate recovery. This typically involves a comprehensive diagnostic workup to accurately stage the cancer, including imaging studies such as computed tomography (CT) or magnetic resonance imaging (MRI) of the neck and chest, and often a positron emission tomography (PET) scan to detect distant metastases. A panendoscopy (laryngoscopy, esophagoscopy, bronchoscopy) with biopsies is performed to delineate tumor extent and rule out synchronous primary tumors.

Beyond oncologic staging, a detailed medical evaluation is essential. This includes a full cardiopulmonary assessment, often with an electrocardiogram (ECG), chest X-ray, and pulmonary function tests, to ensure the patient can tolerate general anesthesia and the physiological changes associated with a permanent tracheostoma. Nutritional status is assessed, and optimization with dietary counseling or pre-operative nutritional support (e.g., oral supplements, nasogastric tube feeding) may be necessary. Patients receive extensive pre-operative counseling regarding the surgical procedure, potential complications, and the profound impact on voice and swallowing. This includes discussions with speech-language pathologists about various voice rehabilitation options (electrolarynx, esophageal speech, tracheoesophageal voice prosthesis) and with stoma nurses regarding tracheostoma care. Dental evaluation and necessary extractions are often performed to minimize post-operative infection risks. Patients are instructed on NPO (nil per os) status before surgery, medication adjustments (e.g., anticoagulants), and typically receive prophylactic broad-spectrum antibiotics. Finally, comprehensive informed consent, detailing all risks and benefits, is obtained.

Contraindications to laryngectomy can be absolute or relative, depending on the patient's overall health, the extent of the disease, and the potential for a curative outcome.

\begin{itemize}
  \item \textbf{Absolute Contraindications:}
    \begin{itemize}
      \item Distant metastatic disease that is widespread and precludes any possibility of cure, as the surgery would not offer a survival benefit.
      \item Unresectable local invasion of the tumor into critical structures such as the carotid artery, prevertebral fascia, or extensive skull base involvement, making complete tumor removal impossible.
      \item Severe, uncorrectable medical comorbidities (e.g., end-stage cardiac failure, severe uncontrolled pulmonary disease) that render the patient unable to tolerate general anesthesia or the physiological stress of a major head and neck surgery.
    \end{itemize}

  \item \textbf{Relative Contraindications and Precautions:}
    \begin{itemize}
      \item Poor functional status or extremely poor performance status (e.g., ECOG performance status 3-4), where the morbidity of surgery outweighs potential benefits.
      \item Significant malnutrition or cachexia, which increases the risk of wound healing complications, infection, and prolonged recovery. Nutritional optimization is crucial.
      \item Severe psychiatric illness or cognitive impairment that prevents the patient from understanding the implications of the surgery or participating in post-operative rehabilitation.
      \item Patient refusal of a permanent tracheostomy or the associated lifestyle changes, despite thorough counseling.
      \item Prior extensive radiation therapy to the neck, which can compromise tissue vascularity and increase the risk of wound breakdown and pharyngocutaneous fistula. Careful flap planning is essential.
      \item Active systemic infection or uncontrolled diabetes, which must be managed prior to surgery to reduce complication rates.
    \end{itemize}
\end{itemize}

\section*{Surgical Technique \& Procedure}
Laryngectomy is performed under general anesthesia, typically with a reinforced endotracheal tube, with the patient positioned supine and the neck extended by placing a shoulder roll to provide optimal surgical exposure. The procedure usually begins with a transverse apron-shaped incision in the lower neck, which may be extended superiorly or inferiorly depending on the extent of the disease and the need for concomitant neck dissection. If indicated, a comprehensive or selective neck dissection is performed to remove lymph nodes that may harbor metastatic cancer, often prior to the laryngeal resection itself, carefully preserving vital structures such as the carotid artery, internal jugular vein, and vagus nerve.

The strap muscles of the neck (sternohyoid, sternothyroid, omohyoid) are then divided or retracted to expose the thyroid gland and the laryngeal cartilages (thyroid and cricoid). The superior laryngeal neurovascular bundles are ligated and divided. The pharynx is carefully entered, typically through the vallecula or piriform sinus, and the larynx is meticulously dissected free from its surrounding attachments, including the hyoid bone, pharyngeal constrictor muscles, and trachea. The trachea is then transected below the cricoid cartilage, usually at the second or third tracheal ring, and the entire larynx is removed en bloc. The remaining pharyngeal defect is then meticulously closed in multiple layers, often with a primary T-closure or linear closure, but sometimes requiring a regional flap (e.g., pectoralis major, radial forearm free flap) for larger defects or in previously irradiated fields to ensure a watertight seal and prevent fistula formation. Finally, the distal trachea is brought out through a separate opening created in the lower neck skin and sutured to the skin edges to form a permanent tracheostoma. A tracheoesophageal puncture (TEP) may be created at this stage if immediate voice prosthesis placement is planned. Surgical drains are placed in the neck wound, and the skin incision is closed.

\section*{Complications \& Risks}
Laryngectomy is a major surgical procedure with a range of potential risks, which can be broadly categorized into general surgical risks and procedure-specific complications.

\begin{itemize}
  \item \textbf{General Surgical Risks:}
    \begin{itemize}
      \item \textbf{Bleeding:} Post-operative hematoma formation requiring re-exploration, or significant blood loss during surgery.
      \item \textbf{Infection:} Wound infection, pneumonia (especially aspiration pneumonia in the early post-operative period), or systemic sepsis.
      \item \textbf{Anesthesia Complications:} Adverse reactions to anesthetic agents, cardiac events (e.g., myocardial infarction, arrhythmias), stroke, deep vein thrombosis, or pulmonary embolism.
    \end{itemize}

  \item \textbf{Procedure-Specific Risks:}
    \begin{itemize}
      \item \textbf{Pharyngocutaneous Fistula:} This is one of the most common and significant complications, occurring when the pharyngeal closure breaks down, leading to leakage of saliva and pharyngeal contents through the neck wound. It prolongs hospitalization and increases infection risk.
      \item \textbf{Pharyngeal/Esophageal Stricture:} Scarring and narrowing of the reconstructed pharynx or esophagus, leading to long-term swallowing difficulties (dysphagia) and requiring dilation or further surgical intervention.
      \item \textbf{Tracheostoma Complications:} Stenosis (narrowing) of the stoma, infection around the stoma, or difficulty with tracheostomy tube placement and care.
      \item \textbf{Voice Prosthesis Issues:} Leakage around or through the tracheoesophageal voice prosthesis, dislodgement, or infection, requiring replacement or revision.
      \item \textbf{Nerve Injury:} Damage to cranial nerves during dissection, such as the hypoglossal nerve (tongue weakness), or spinal accessory nerve (shoulder weakness if neck dissection is performed).
      \item \textbf{Chyle Leak:} If a neck dissection is performed, injury to the thoracic duct can lead to leakage of lymphatic fluid (chyle) into the neck wound, requiring drainage or re-exploration.
      \item \textbf{Wound Dehiscence:} Partial or complete breakdown of the surgical incision, particularly in previously irradiated tissues.
      \item \textbf{Hypothyroidism:} If the thyroid gland is partially or completely removed during the procedure, patients will require lifelong thyroid hormone replacement therapy.
      \item \textbf{Cosmetic Deformity:} Changes in neck contour and appearance due to tissue removal and scar formation.
    \end{itemize}
\end{itemize}

\section*{Postoperative Care \& Recovery}
Immediately following laryngectomy, patients are transferred to the Post-Anesthesia Care Unit (PACU) for close monitoring of vital signs, airway patency through the new tracheostoma, and assessment for immediate post-operative complications like bleeding. A tracheostomy tube is in place to maintain the airway, and pain is managed with intravenous analgesics. Patients will have surgical drains in their neck and a nasogastric tube (NGT) or gastrostomy tube (PEG) for nutritional support, as oral intake is typically withheld for 7-10 days to allow for pharyngeal healing. Early mobilization is encouraged to prevent pulmonary complications and deep vein thrombosis.

Over the first 24-48 hours, monitoring continues for signs of pharyngocutaneous fistula (e.g., excessive drain output, salivary leakage), hematoma, or infection. Stoma care education begins, involving nurses and speech-language pathologists, to teach patients and their families how to clean and manage the tracheostoma. Once the pharyngeal closure is deemed stable, typically after a swallow study (e.g., modified barium swallow or esophagram) confirms no leakage, the NGT can be removed, and a gradual progression to oral diet is initiated. Speech rehabilitation with an electrolarynx usually begins within the first few days, while training for esophageal speech or the use of a tracheoesophageal voice prosthesis (if placed) commences once healing is sufficient. Long-term recovery involves ongoing speech and swallowing therapy, psychosocial support to adjust to changes in communication and body image, and regular follow-up with the surgical and oncology teams for surveillance. Full physical recovery and adaptation to the new lifestyle can take several months to a year.

During laryngectomy, the surgeon meticulously documents the extent of the tumor, its precise location, involvement of adjacent structures, and the presence of any suspicious lymph nodes. This intraoperative assessment guides the extent of resection and any concomitant neck dissection. Following removal, the resected larynx and any associated lymph nodes are sent for histopathological examination. The pathology report is paramount, detailing the exact type of cancer (most commonly squamous cell carcinoma), its histological grade, size, and multifocality. Most critically, the report specifies the status of the surgical margins, indicating whether the tumor was completely removed with a rim of healthy tissue (clear margins) or if tumor cells are present at the edge of the resection (positive or close margins). The report also quantifies the number of involved lymph nodes, their location, and the presence of extranodal extension, which are crucial prognostic indicators and guide decisions regarding adjuvant therapy.

A normal, successful postoperative course following laryngectomy involves a predictable sequence of healing and functional adaptation. Initially, patients will experience pain at the surgical site, managed effectively with analgesics, which gradually subsides over the first few weeks. The neck wound should heal cleanly, and the tracheostoma should be patent, well-formed, and easily manageable by the patient or caregiver. The absence of a pharyngocutaneous fistula is a key indicator of successful pharyngeal healing. Oral intake is typically resumed within 7-14 days after a negative swallow study, with a gradual progression from liquids to soft foods. Patients will learn to manage their tracheostoma, including cleaning, humidification, and suctioning. Voice rehabilitation, whether through an electrolarynx, esophageal speech, or a tracheoesophageal voice prosthesis, progresses over weeks to months, with the goal of achieving functional and intelligible communication. While some degree of swallowing difficulty may persist, a successful course allows for adequate oral nutrition and a return to most daily activities with appropriate support and adaptation.

Post-laryngectomy follow-up is a critical, long-term process designed to monitor for cancer recurrence, manage complications, and optimize functional outcomes.

\begin{itemize}
  \item \textbf{Regular Clinic Visits:} Frequent visits with the head and neck surgeon and oncologist, initially weekly or bi-weekly, then gradually spaced out to monthly, quarterly, and eventually annually for several years.
  \item \textbf{Suture/Staple Removal:} Typically performed within 1-2 weeks post-operatively.
  \item \textbf{Speech and Swallowing Therapy:} Ongoing sessions with a speech-language pathologist to refine voice production techniques, improve swallowing efficiency, and manage any dysphagia or aspiration issues.
  \textbf{Nutritional Counseling:} To ensure adequate caloric intake and address any weight loss or dietary challenges.
  \item \textbf{Periodic Imaging:} Surveillance imaging, such as CT or MRI of the neck and chest, usually performed every 3-6 months for the first few years, then annually, to detect any local or regional recurrence or distant metastases.
  \item \textbf{Fiberoptic Laryngoscopy/Pharyngoscopy:} Regular endoscopic examination of the pharynx and stoma to check for mucosal changes or early signs of recurrence.
  \item \textbf{Tracheoesophageal Puncture (TEP) Management:} Regular maintenance, cleaning, and replacement of the voice prosthesis as needed, typically every 3-6 months.
  \item \textbf{Psychosocial Support:} Access to support groups, counseling, or psychological services to help patients cope with the significant physical and emotional changes.
  \item \textbf{Activity Resumption:} Gradual return to normal activities, with specific guidance on exercise, swimming (with stoma protection), and avoiding activities that could compromise the stoma.
\end{itemize}
Patients are educated to report any new symptoms immediately, such as persistent pain, bleeding from the stoma or mouth, difficulty breathing, changes in swallowing, unexplained weight loss, or new lumps in the neck.

\section*{Outcomes \& Prognosis}
Laryngeal cancer, for which laryngectomy is a primary treatment, represents a significant portion of head and neck malignancies. Globally, it accounts for approximately 1-2\% of all cancers, with an estimated 180,000 new cases diagnosed annually. In the United States, there are roughly 12,000-13,000 new cases each year. The incidence of laryngeal cancer shows a strong male predominance, with men being affected 4 to 5 times more frequently than women, although this gap has been narrowing in recent decades. The disease typically affects individuals in their 50s, 60s, and 70s. The primary risk factors are chronic tobacco smoking and heavy alcohol consumption, which act synergistically to increase risk. Human papillomavirus (HPV) infection, particularly HPV-16, is an emerging risk factor, especially for certain subsites like the glottis and supraglottis, though its role in laryngeal cancer is less pronounced than in oropharyngeal cancer. The mortality rate for laryngeal cancer has been declining due to earlier detection and improved treatment modalities, but advanced-stage disease still carries a substantial risk of death. Laryngectomy is performed in a significant proportion of these advanced cases, contributing to the overall morbidity profile of head and neck cancer treatment.

The outcomes and prognosis following laryngectomy are highly dependent on the stage of the laryngeal cancer at diagnosis, the presence of lymph node metastases, the status of surgical margins, and the patient's overall health. For advanced-stage laryngeal cancer (T3 or T4) treated with total laryngectomy, 5-year survival rates typically range from 60\% to 80\%. The presence of positive surgical margins or extranodal extension in resected lymph nodes significantly worsens the prognosis and often necessitates adjuvant radiation therapy or chemoradiation. Functional outcomes are a critical aspect of prognosis, with successful voice rehabilitation (e.g., intelligible tracheoesophageal speech) achieved in a high percentage of patients, though some may struggle with communication. Swallowing function can be permanently altered, with a subset of patients experiencing chronic dysphagia or requiring long-term feeding tube dependence, particularly after extensive resections or prior radiation. Recurrence rates vary but are generally lower for patients with clear margins and negative nodes. Prognostic factors include tumor stage, nodal status, HPV status (for certain subsites), and the patient's ability to complete adjuvant therapy. While laryngectomy offers a high chance of local disease control for advanced laryngeal cancer, it comes with significant functional morbidity that requires comprehensive, lifelong rehabilitation and support.

\section*{References}
\begin{enumerate}
  \item MedlinePlus. Laryngectomy. U.S. National Library of Medicine; 2024. Available from: \url{https://medlineplus.gov/ency/article/002932.htm}
  \item Cummings CW, Flint PW, Haughey BH, et al., editors. Cummings Otolaryngology: Head \& Neck Surgery. 7th ed. Elsevier; 2021.
  \item National Comprehensive Cancer Network (NCCN). NCCN Clinical Practice Guidelines in Oncology: Head and Neck Cancers. Version 2.2024.
  \item Singer MI, Blom ED. An endoscopic technique for restoration of voice after laryngectomy. Ann Otol Rhinol Laryngol. 1980;89(6 Pt 1):529-33.
\end{enumerate}

\clearpage
